\documentclass[12pt,a4paper]{article}

% PACOTES
\usepackage[portuguese]{babel}
\usepackage[utf8]{inputenc}
\usepackage[T1]{fontenc}
\usepackage{makeidx}
\usepackage{xspace}
\usepackage{graphicx,color,times}
\usepackage{fancyhdr}
\usepackage{latexsym}
\usepackage[printonlyused]{acronym} % Pacote acronym aqui
\usepackage{float}
\usepackage{listings}
% remove duplicate "index"
\usepackage[nottoc]{tocbibind}
\usepackage{natbib}
\usepackage{hyperref}
\usepackage{xcolor}
\usepackage{fix-cm}
\usepackage{fourier}
\usepackage[scaled=.92]{helvet}
% show references names
\usepackage{nameref}

% CONFIGURAÇÕES HYPERREF
\hypersetup{
    colorlinks=true,            
    linkcolor=blue,            
    citecolor=blue,             
    urlcolor=blue,              
    filecolor=blue,             
    breaklinks=true,            
    pdftitle={A Governancia da Internet},   
    pdfauthor={Tomás Simões Duarte},           
    pdfsubject=
    {
        As decisões no design e desenvolvimento 
        de tecnologias que são necessárias para manter a Internet operacional 
        com a implementação
        de políticas novas em volta das mesmas
    }, 
    pdfkeywords={Governancia, Internet},   
}

% DEFINIÇÃO DE ACRÔNIMOS
\begin{document}

% Frontmatter
\thispagestyle{empty}
\setcounter{page}{-1}

\begin{center}
  \begin{Huge}
    \textbf{Universidade da Beira Interior}
  \end{Huge}
\end{center}

\begin{center}
  \begin{Huge}
    Departamento de Informática
  \end{Huge}
\end{center}

\vspace{0,07cm}
\begin{figure}[!htb]
  \centering
  \includegraphics[width=191pt]{images/ubi-fe-di.png}
\end{figure}

\vspace{0.5cm}
\begin{center}
  \begin{Large}
    \textbf{A Governancia da Internet}
  \end{Large}
\end{center}


\vspace{0.5cm}
\begin{center}
  \begin{normalsize}
    \begin{large}
      Elaborado por:
    \end{large}
  \end{normalsize}
\end{center}

\vspace{0.2cm}
\begin{center}
  \begin{large}
    \textbf{Tomás Simões [52585]}
  \end{large}
\end{center}

\vspace{0,5cm}
\begin{center}
  \begin{normalsize}
    \begin{large}
      Orientador:
    \end{large}
  \end{normalsize}
\end{center}

\vspace{0.2cm}
\begin{center}
  \begin{large}
    \textbf{Professor Doutor Bruno Miguel Silva}
  \end{large}
\end{center}



\vspace{0.5cm}
\begin{center}
  \begin{normalsize}
    \today
  \end{normalsize}
\end{center}

\let\cleardoublepage\clearpage
\tableofcontents

% Mainmatter
\acresetall
\section{Introdução}
\label{chap:intro}
Este relatório, com o tema de "Governance da Internet" foi feito no contexto da unidade curricular de Redes de Computador, no projeto
enquanto frequentava a \ac{UBI}.

Atualmente vivemos em uma sociedade que é governada pela tecnologica, apesar de a maioria dos individuos
não saberem quem a governa.

Ao longo dos anos tivemos diferentes perpetivas sobre o uso desta,
desde o uso militar até simplemente podermos usar o Youtube num frigorifico. Por essa razão
é necessário atualizar e dialogar sobre como deviamos nos comportar face a questões políticas, sociais e filosoficas da rede de redes
de modo a termos um espaço seguro para estar.

Este documento fala sobre o aparecimento do termo "Governancia da Internet", um campo de estudo e pesquisa,
complexo e com diversas intrepetações sobre quem  dentro da Internet.

\subsection{Organização do Documento}
\label{sec:organ}

De modo a refletir o trabalho que foi feito, este documento encontra-se estruturado da seguinte forma:
\begin{enumerate}
  \item O 1º capítulo -- \textbf{Introdução} -- apresenta o projeto, a motivação para a sua escolha, o enquadramento para o mesmo, os seus objetivos e a respetiva organização do documento.
  \item O 2º capítulo -- \textbf{Tecnologias Utilizadas} -- descreve os conceitos mais importantes no âmbito deste projeto, bem como as tecnologias utilizadas durante do desenvolvimento da aplicação.
  \item ...
\end{enumerate}

\section{O que é a Governancia da Internet?}
\label{sec:internetgovernance}
A Governancia da Internet é definida como "decisões no \textit{design} e desenvolvimento 
de tecnologias que são necessárias para manter a Internet operacional com a implementação
de políticas novas em volta das mesmas".

Ao longo do tempo, pesquisadores assim como profissionais na área têm tentado expandir o 
conhecimento deste extenso dominio, apresentando divisões para os problemas em questão.

Uma dessas classificações, proposta por DeNardis e Raymond, segmenta a Governança da Internet em seis áreas funcionais:

\begin{enumerate}
  \item Controle de Recursos Críticos da Internet (CIRs);
  \item Estabelecimento de \textit{standards} da Internet;
  \item Coordenação de Acesso e Interligação;
  \item Governança de Cibersegurança;
  \item Intermediação de Informações;
  \item Direitos de Propriedade Intelectual baseados em Arquitetura;
\end{enumerate}

Estas áreas têm diferentes estratégias e modelos de admistração, cada um adaptado
para ser mais eficiente na sua área especifica. No entanto, todos os pontos 
convergem para um objetivo comum.

Examinando a trajetória da Internet ao longo do tempo, percebe-se que ela foi 
observada em diferentes contextos históricos e culturais, acompanhada de 
diversas expectativas e decepções.

Por conta disso, foi preciso mantar uma abordagem de flexivel face ao desenvolvimento
da arquitetura da Internet, permitindo que novos atores com diferentes ideias entrassem 
e modificassem a estrutura conforme desejado.





\section{O inicio da Internet}
\label{sec:beggining}
\section{Evolução da Governancia da Internet}
\label{sec:governanceevo}
O termo "Governancia da Internet" não está em nenhum artigo escolar até 1995. No entanto começaram a surgir pesquisadores com propostas em como aplicar 
governancia e leis na Internet.

O surgimento desta área de pesquisa converge para acontecimentos em meados do século 19, 
em que a Internet começou a emergir como uma massiva mídia
causada por 3 principais eventos:

\begin{enumerate}
  \item O desenvolvimento e adoção do protocologo WWW (1989-1993);
  \item[]Permitiu a criação de uma rede de informações interligadas que poderiam
  ser acessadas por meio de navegadores da \textit{web}.
  
  \item A publicação de \textit{web browsers} gratuitos para utilização (1991);
  
  \item[]Os \textit{web browsers} geralmente eram comercializados e a disponibilidade dos mesmos 
  facilitou o acesso à Internet e abriu portas para novos individuos na sua população.
  \item Privatização da Internet e a sua abertura para uso comercial pela \textit{US National Science Foundation} (1995) 

  \item[]Antigamente, a infraestrutura principal da Internet era predominanmente controlada e financiada pelo
  governo ou instituições acadêmicas, mas a decisão de permitir o uso comercial abriu caminho
  para o desenvolvimento do comércio eletrónico.


\end{enumerate}

% \include{chap-5}
% \include{chap-6}
% \chapter*{Conclusão}
\addcontentsline{toc}{chapter}{Conclusão}
\label{chap:conc-trab-futuro}

Esta secção contém a resposta à questão: \\
\emph{Quais foram as conclusões princípais a que o(a) aluno(a) chegou no fim deste trabalho?}

Esta secção responde a questões como:\\
\emph{O que é que ficou por fazer, e porque?}\\
\emph{O que é que seria interessante fazer, mas não foi feito por não ser exatamente o objetivo deste trabalho?}\\
\emph{Em que outros casos ou situações ou cenários -- que não foram estudados no contexto deste projeto por não ser seu objetivo -- é que o trabalho aqui descrito pode ter aplicações interessantes e porque?}
teste 1 2 \cite{musiani2013network}.
\bibliography{references}
\end{document}